\section{VOA de réseau, involution canonique et branche Moonshine d'ordre deux}
\label{sec:voa-moonshine-orden-dos-rev11}
\index{VOA!réseau de Leech}
\index{Moonshine!classe 2A}

La branche d'incidence construit d'abord le réseau de Leech
\(\Lambda_{24}\). À partir de cette donnée, la théorie classique des algèbres
d'opérateurs de vertex fournit un prolongement canonique. Si
\(\mathfrak h=\Lambda_{24}\otimes_{\mathbb Z}\mathbb C\) et \(M(1)\) est
le module de Fock de son algèbre de Heisenberg, le VOA de réseau est
\begin{equation}
 V_{\Lambda}=M(1)\otimes\mathbb C[\Lambda_{24}],
 \label{eq:voa-reticular-leech-rev11}
\end{equation}
de charge centrale \(24\). Pour \(\alpha\in\Lambda_{24}\), ses opérateurs
de vertex s'écrivent
\[
 Y(e^\alpha,z)=E^-(-\alpha,z)E^+(-\alpha,z)e^\alpha z^{\alpha_0}.
\]
L'absence de racines dans \(\Lambda_{24}\) implique l'annulation du secteur
de poids un dans l'orbifold de Moonshine.

L'involution
\[
 \theta:\Lambda_{24}\longrightarrow\Lambda_{24},
 \qquad \theta(\lambda)=-\lambda,
\]
est intrinsèque et se prolonge à \(V_{\Lambda}\). Sa trace graduée s'obtient
par comptage direct des modes oscillateurs.

\begin{theorem}[Trace de l'involution de négation]
\label{thm:traza-theta-leech-rev11}
\begin{equation}
 t_2(\tau)
 :=\operatorname{Tr}_{V_\Lambda}
 \bigl(\theta q^{L_0-1}\bigr)
 =q^{-1}\prod_{n\geq1}(1+q^n)^{-24}.
 \label{eq:t2-producto-leech-rev11}
\end{equation}
\end{theorem}
\begin{proof}
L'involution apparie \(e^\lambda\) et \(e^{-\lambda}\) ; comme le réseau
est sans torsion, seul \(\lambda=0\) contribue à la trace du secteur
réticulaire. Chacun des vingt-quatre oscillateurs d'énergie \(n\) apporte
\(\sum_{k\geq0}(-1)^kq^{nk}=(1+q^n)^{-1}\). Le décalage du vide
produit le facteur \(q^{-1}\).
\end{proof}

La complétion de Fricke associée à cette trace est
\begin{equation}
 T_{2A}(\tau)=t_2(\tau)+24+\frac{2^{12}}{t_2(\tau)}
 =q^{-1}+4372q+96256q^2+1240002q^3+\cdots.
 \label{eq:T2A-hmt-rev11}
\end{equation}
Le terme \(24\) annule le terme constant de \(t_2\) ; le facteur
\(2^{12}\) est la dégénérescence du secteur tordu de vingt-quatre bosons.
La même série de bord détermine l'invariant modulaire au moyen de
\begin{equation}
 \boxed{j(\tau)=\frac{(t_2(\tau)+256)^3}{t_2(\tau)^2}}.
 \label{eq:j-desde-t2-rev11}
\end{equation}
Les équations \eqref{eq:t2-producto-leech-rev11}--
\eqref{eq:j-desde-t2-rev11} construisent la branche canonique d'ordre deux sans
paramètres continus.

L'orbifold de Frenkel--Lepowsky--Meurman est défini par
\[
 V^\natural=(V_\Lambda)^+\oplus(V_\Lambda^T)^+.
\]
C'est un VOA holomorphe de charge centrale \(24\), qui satisfait
\((V^\natural)_1=0\), possède pour caractère
\[
 \operatorname{ch}_{V^\natural}(\tau)=j(\tau)-744
\]
et son groupe d'automorphismes est le groupe Monstre
\cite{FLM1988,ConwayNorton1979,Borcherds1992}. Dans cette chaîne, HMT fixe le
réseau de Leech et l'involution \(\theta\) ; les théorèmes classiques de VOA et
de FLM identifient le prolongement résultant à \(V^\natural\) et à son groupe
d'automorphismes. L'énoncé n'exige pas d'action du Monstre sur les
chiffres : l'évaluation archimédienne et la branche exceptionnelle descendent d'un
antécédent HMT commun et conservent des invariants différents.

La fonction voisine
\(T_{3A}=t_3+12+729/t_3\) appartient à la théorie classique des Hauptmoduln
d'ordre trois. La réalisation HMT active dans le corridor tritique est la classe
\(3B\) ; une identification à \(3A\) exige de déclarer le relèvement spécifique
qui change de classe.
